\documentclass[11pt]{article}
\usepackage{amsmath,amssymb,hyperref}
\begin{document}
\title{Exact pure anticoherent states at $N=27,29,30$}\author{}\date{}\maketitle
Throughout, $D_N^r$ denotes the normalized Dicke vector with $r$ excitations.
Anticoherence of order $t$ means $\rho_t=I_{\mathrm{sym},t}/(t+1)$.
The constructions below establish exact attainment. A separate universal
exclusion theorem is needed whenever a global maximum order is claimed.

\section*{Two fifth-order constructions}
Define
\begin{align*}
w_{27}&=\frac{52-i\sqrt{3433}}{19\sqrt{17}},\\
|\psi_{27}\rangle&=\frac{\sqrt{327}(D_{27}^1-D_{27}^{26})
-i\sqrt{494}(D_{27}^6-D_{27}^{21})
-iw_{27}\sqrt{429}(D_{27}^{11}-D_{27}^{16})}{50},\\
w_{29}&=\frac{3094-i\sqrt{2170247}}{\sqrt{11743083}},\\
|\psi_{29}\rangle&=\frac{\sqrt{3598}(D_{29}^2-D_{29}^{27})
-i\sqrt{2081}(D_{29}^7-D_{29}^{22})
-4iw_{29}\sqrt{231}(D_{29}^{12}-D_{29}^{17})}{25\sqrt{30}}.
\end{align*}
Both phase parameters have modulus one. Their probabilities are respectively
\[
\frac{(327,494,429,429,494,327)}{2500},\qquad
\frac{(3598,2081,3696,3696,2081,3598)}{18750}.
\]
They are normalized and symmetric about $r=N/2$. The diagonal moment
identities are
\[
\sum_r |c_r|^2 (r)_q=\frac{(N)_q}{q+1},\qquad 0\le q\le5,
\]
where $(r)_q$ is the falling factorial. These are equivalent to the six
uniform diagonal entries of $\rho_5$. Excitation indices differ by multiples
of five, so its only possible off-diagonal entries are $(0,5)$ and $(5,0)$.
Substituting the displayed phases into
\[
(\rho_5)_{05}=\sum_{r=0}^{N-5}\binom{N-5}{r}
\frac{c_r\overline{c_{r+5}}}{\sqrt{\binom Nr\binom N{r+5}}}
\]
gives zero exactly. Thus both states are fifth-order anticoherent.
They are not sixth-order:
\[
(\rho_6^{27})_{00}-\frac17=\frac{17}{1593900},\qquad
(\rho_6^{29})_{00}-\frac17=-\frac{578}{1781325}.
\]
The projective Majorana multiplicities are $1^{27}$ and $(2^2,1^{25})$.
The independent verifier proves simplicity of the remaining roots using a
reciprocal quintic in $z^5$, with nonzero quadratic discriminant and no
branch point at $z^5=\pm1$.

\section*{A seventh-order octahedral construction}
Set
\[
F_6(z)=z(1-z^4),\quad F_8(z)=1+14z^4+z^8,\quad
\alpha=\frac{\sqrt{326271}}4,\quad
\beta=\frac{-118635+14355i\sqrt{287}}{256\sqrt{326271}}.
\]
Write $P(z)=\alpha F_6(z)^5+\beta F_6(z)F_8(z)^3$
and define $c_r=[z^r]P(z)/\sqrt{\binom{30}r}$.
Equivalently,
\begin{align*}
|\psi_{30}\rangle={}&\frac{\beta}{\sqrt{30}}(D_{30}^1-D_{30}^{29})
+\frac{\alpha+41\beta}{\sqrt{142506}}(D_{30}^5-D_{30}^{25})\\
&+\frac{-5\alpha+549\beta}{\sqrt{14307150}}(D_{30}^9-D_{30}^{21})
+\frac{10\alpha+2237\beta}{\sqrt{119759850}}(D_{30}^{13}-D_{30}^{17}).
\end{align*}
The polynomial definition is sufficient without any group-theoretic
assumption: direct exact substitution into the partial-trace formula below
proves normalization and $\rho_7=I_8/8$. In deriving the coefficients,
the octahedral binary forms leave two moment conditions after normalization;
the relevant Gram entries are
\[
|\alpha|^2=\frac{326271}{16},\quad
|\beta|^2=\frac{14025}{4096},\quad
\operatorname{Re}(\alpha\overline\beta)=-\frac{118635}{1024}.
\]
Their determinant is $59140949175/1048576>0$, giving the stated nonreal
relative phase. Moreover
\[
(\rho_8)_{00}-\frac19=-\frac{11}{538200}\ne0.
\]
Thus the order is exactly seven. All30 projective Majorana roots are simple.
After extracting $F_6$, the remaining degree-six polynomial in $q=z^4$
is coprime to its derivative over $\mathbb Q(i\sqrt{287})$ and is nonzero
at $q=0,1$.

\section*{Independent reproduction}
For any normalized Dicke vector the reduction is
\[
(\rho_q)_{ab}=\sum_{\ell=0}^{N-q}c_{a+\ell}\overline{c_{b+\ell}}
\sqrt{\frac{\binom qa\binom qb\binom{N-q}{\ell}^2}
{\binom N{a+\ell}\binom N{b+\ell}}}.
\]
\texttt{build\_witnesses.py} writes the explicit radical vectors and checks
Schmidt matrices. Independently, \texttt{verify\_witnesses.py} reads the saved
vectors, evaluates the direct formula for orders $t,t+1$ using exact SymPy
arithmetic and verifies the stated root multiplicities. Lower orders follow
by partial trace. Decimal coefficients and sphere coordinates are renderings
of these exact vectors and are not proof premises.
\end{document}
